\begin{theorem}[Cyclic-prime norm filtration theorem]
\label{mod:ramification-normfiltration}
Let $K/F$ be a prime cyclic extension of nonarchimedean local fields,
let $t\in\mathbf N$ be a supplied lower-break witness, and assume explicitly
that
\[
 f(K/F)=1.
\]
Fix an integral uniformizer $\pi_K\in\mathcal O_K$ satisfying
\[
 \operatorname{Algebra.adjoin}_{\mathcal O_F}\{\pi_K\}=\mathcal O_K.
\]
Put
\[
 \ell=[K:F],\qquad
 D=(\ell-1)(t+1),\qquad
 \Psi(r)=\Psi_{t,\ell}(r).
\]
For $a\le b$, write
\[
 Q_E(a,b)=U_E^a/U_E^b.
\]
All maps below are induced by the field norm.  Their numerator and
denominator containments are supplied separately, so every quotient map is
well-defined and independent of representatives.

Lean first identifies the two filtration-containment predicates used by
modules~\ref{mod:ramification-normbelowbreak} and
\ref{mod:ramification-normabovebreak}.  For every natural target depth $r$ it
then proves
\[
 N_{K/F}(U_K^{\Psi(r)})\subseteq U_F^r,
 \qquad
 N_{K/F}(U_K^{\Psi(r)+1})\subseteq U_F^{r+1}.
\]
Thus norm defines the all-depth graded map
\[
 \overline N_r^{\mathrm H}:
 Q_K(\Psi(r),\Psi(r)+1)\longrightarrow
 Q_F(r,r+1)=\operatorname{gr}^rU_F,
 \qquad [u]\longmapsto[N_{K/F}(u)].
\]
This single Herbrand-indexed map is connected explicitly to all three
case-specific maps.  It is bijective for $r<t$ and for $t<r$.  At $r=t$
its range is exactly the range of the canonical critical map, and its
kernel and cokernel both have order $\ell$.  Thus no dependent transport
from the identity $\Psi(r)=r$ is left implicit in the assembled API.
More generally, for $r\le s$ it defines
\[
 N^{\mathrm H}_{r,s}:
 Q_K(\Psi(r),\Psi(s))\longrightarrow Q_F(r,s),
 \qquad [u]\longmapsto[N_{K/F}(u)].
\]
These maps commute with every canonical projection obtained by deepening
the source and target denominators.  Explicitly, if
$r\le s_1\le s_2$, then
\[
 \operatorname{pr}_F\circ N^{\mathrm H}_{r,s_2}
 =
 N^{\mathrm H}_{r,s_1}\circ\operatorname{pr}_K.
\]
The same projection identity is exported for arbitrary certified
heterogeneous numerator and denominator pairs, both in the below-break
containment notation and in the above-break containment notation.  On the
successive target quotient,
\[
 N^{\mathrm H}_{r,r+1}
 =
 \overline N_r^{\mathrm H}\circ
 \operatorname{pr}:
 Q_K(\Psi(r),\Psi(r+1))
 \longrightarrow Q_K(\Psi(r),\Psi(r)+1)
 \longrightarrow\operatorname{gr}^rU_F.
\]

Every natural depth belongs to exactly one of the three cases
\[
 r<t,\qquad r=t,\qquad t<r.
\]
The inverse-Herbrand boundary formulas and denominator comparisons used in
this case split are
\[
 \Psi(t)=t,\qquad
 \Psi(t+1)=t+\ell,
\]
\[
 \Psi(r)=t+\ell(r-t)\qquad(t\le r),
\]
and
\[
 \Psi(r)+1\le\Psi(r+1)\qquad(r\ge0).
\]
For $t\le r$, primality of $\ell$ gives the strict inequality
\[
 \Psi(r)+1<\Psi(r+1).
\]
In particular, the denominator of the graded norm above the break is
$\Psi(r)+1$, never the next Herbrand depth $\Psi(r+1)$.

The different and trace conclusions are
\[
 d_{K/F}=D=(\ell-1)(t+1)
\]
and, for every $a\in\mathbf Z$,
\[
 \operatorname{Tr}_{K/F}(\mathfrak p_K^a)
 =
 \mathfrak p_F^{\left\lfloor(a+D)/\ell\right\rfloor}.
\]
The trace statement is an equality of mapped lattices.  The replacement of
the ramification index by $\ell$ uses the explicit equality $f(K/F)=1$.

Strictly below the break, for every $i<t$, the exact same-numbered graded
norm
\[
 \overline N_i:
 \operatorname{gr}^iU_K\longrightarrow\operatorname{gr}^iU_F
\]
is bijective.  Splicing exactly the indices $i=r,\ldots,t-1$ proves that
\[
 Q_K(r,t)\longrightarrow Q_F(r,t)
\]
is bijective for every $0\le r\le t$, including the trivial endpoint
$r=t$.  Norm also induces a bijection
\[
 K^\times/U_K^t\longrightarrow F^\times/U_F^t.
\]
When $t=0$, this is the valuation quotient; it does not assert that the
critical residue-unit norm is bijective.

Lean additionally describes the honest full norm image on every layer
$r\le t$.  Define
\[
 M_{r,t+1}:Q_K(r,t+1)\longrightarrow Q_F(r,t+1)
\]
to be the norm-induced quotient map, and let
\[
 J_r=
 \left\{
 u\in U_F^r:
 [u]\in\operatorname{im}(M_{r,t+1})
 \right\}
 \le U_F^r.
\]
Then the range of the restricted norm homomorphism
$U_K^r\to U_F^r$ is exactly $J_r$:
\[
 \operatorname{im}\bigl(U_K^r\longrightarrow U_F^r\bigr)=J_r
 \qquad(r\le t).
\]
Projection of $M_{r,t+1}$ from denominator $t+1$ to denominator $t$
recovers the preceding spliced map
$Q_K(r,t)\to Q_F(r,t)$.

At the critical depth, put
\[
 R_t=\operatorname{im}\left(
 \overline N_t:
 U_K^t/U_K^{t+1}\longrightarrow U_F^t/U_F^{t+1}
 \right)
\]
and define
\[
 J_t^{\mathrm{crit}}
 =
 \left\{
 u\in U_F^t:[u]\in R_t
 \right\}
 \le U_F^t.
\]
Let $I_t\le F^\times$ be the image of
$J_t^{\mathrm{crit}}\hookrightarrow F^\times$.  Lean proves
\[
 N_{K/F}(U_K^t)=I_t,
\]
equivalently,
\[
 \operatorname{im}\bigl(U_K^t\longrightarrow U_F^t\bigr)
 =J_t^{\mathrm{crit}}.
\]
The all-below-layers definition agrees with this definition at the endpoint:
\[
 J_t=J_t^{\mathrm{crit}}.
\]
The reverse inclusion in the critical full-image equality uses the exact
successor image
\[
 \boxed{\,N_{K/F}(U_K^{t+1})=U_F^{t+1}\,}.
\]
This equality is also exported in elementwise existence and if-and-only-if
forms and includes both $t=0$, where it is the $U^1$ endpoint, and every
positive break.

The critical graded kernel is exactly the ramification image:
\[
 \ker(\overline N_t)
 =
 \operatorname{im}\left(
 G_t/G_{t+1}\longrightarrow\operatorname{gr}^tU_K
 \right).
\]
The ramification map is injective, and its composition with
$\overline N_t$ is exact at the middle term.  Moreover,
\[
 \#\ker(\overline N_t)=\ell,
 \qquad
 \#\left(
 \operatorname{gr}^tU_F/R_t
 \right)=\ell.
\]
The cokernel is therefore the target quotient by the actual norm range,
not a quotient in the reverse direction.

For every $r\le t$, the full below-break image is further characterized by
the subgroup equalities
\[
 N_{K/F}(U_K^r)\cap U_F^t=I_t,
 \qquad
 N_{K/F}(U_K^r)\vee U_F^t=U_F^r,
\]
where $\vee$ denotes subgroup join in $F^\times$.  These formulas give the
exact full image without asserting the false equality
$N_{K/F}(U_K^r)=U_F^r$.  The corresponding full-field equalities are
\[
 N_{K/F}(K^\times)\cap U_F^t=I_t,
 \qquad
 N_{K/F}(K^\times)\vee U_F^t=F^\times.
\]

Strictly above the break, for every $r>t$, the canonical graded norm is
bijective and has full range:
\[
 Q_K(\Psi(r),\Psi(r)+1)
 \xrightarrow{\ \sim\ }
 \operatorname{gr}^rU_F.
\]
Complete lifting gives the two exact full-image equalities
\[
 \boxed{\,N_{K/F}(U_K^{\Psi(r)})=U_F^r\,},
 \qquad
 \boxed{\,N_{K/F}(U_K^{\Psi(r)+1})=U_F^{r+1}\,}.
\]
At the first strict upper depth $r=t+1$, the first source is
$U_K^{t+\ell}$, whereas the independently proved critical-successor
equality uses $U_K^{t+1}$.  These endpoints are not identified.

Finally, \texttt{CyclicPrimeNormFiltration} is a proposition-valued record
packaging the depth trichotomy, the displayed Herbrand formulas and
containments, the three all-depth graded-map bridges, the different and
trace equalities, all below-break
bijectivity and exact-image statements, the critical kernel, exactness,
kernel and cokernel cardinalities, the critical successor image, both
above-break graded and full-image statements, and the two global field-image
equalities.  The canonical homomorphisms, representative formulas, image
subgroups, and projection compatibilities remain separate public
definitions and the theorem
\texttt{cyclicPrimeNormFiltration} constructs this record from
modules~\ref{mod:ramification-different},
\ref{mod:ramification-traceideals},
\ref{mod:ramification-normbelowbreak}, and
\ref{mod:ramification-normabovebreak}.

\LeanFile{LanglandsFirstMainLemma/Ramification/NormFiltration.lean}
\lean{LanglandsFirstMainLemma.normMapsUnitFiltration_iff_aboveBreak}
\lean{LanglandsFirstMainLemma.normUnitFiltrationQuotient_projection}
\lean{LanglandsFirstMainLemma.cyclicPrimeAboveBreakQuotient_projection}
\lean{LanglandsFirstMainLemma.normMapsUnitFiltration_herbrand}
\lean{LanglandsFirstMainLemma.normMapsUnitFiltration_herbrand_succ}
\lean{LanglandsFirstMainLemma.cyclicPrimeGradedNorm}
\lean{LanglandsFirstMainLemma.cyclicPrimeGradedNorm_mk}
\lean{LanglandsFirstMainLemma.cyclicPrimeGradedNorm_bijective_belowBreak}
\lean{LanglandsFirstMainLemma.cyclicPrimeGradedNorm_range_atBreak}
\lean{LanglandsFirstMainLemma.cyclicPrimeGradedNorm_cokernel_card_atBreak}
\lean{LanglandsFirstMainLemma.cyclicPrimeGradedNorm_kernel_card_atBreak}
\lean{LanglandsFirstMainLemma.cyclicPrimeGradedNorm_bijective_aboveBreak}
\lean{LanglandsFirstMainLemma.cyclicPrimeNormFiltrationQuotient}
\lean{LanglandsFirstMainLemma.cyclicPrimeNormFiltrationQuotient_mk}
\lean{LanglandsFirstMainLemma.cyclicPrimeNormFiltrationQuotient_projection}
\lean{LanglandsFirstMainLemma.herbrandPsiNat_add_one_le_succ}
\lean{LanglandsFirstMainLemma.cyclicPrimeNormFiltrationQuotient_succ}
\lean{LanglandsFirstMainLemma.norm_unitFiltration_surjective_break_succ}
\lean{LanglandsFirstMainLemma.mem_unitFiltration_break_succ_iff_exists_norm}
\lean{LanglandsFirstMainLemma.norm_unitFiltration_map_eq_break_succ}
\lean{LanglandsFirstMainLemma.criticalNormTargetImage}
\lean{LanglandsFirstMainLemma.norm_unitFiltration_range_eq_criticalNormTargetImage}
\lean{LanglandsFirstMainLemma.criticalNormImage}
\lean{LanglandsFirstMainLemma.norm_unitFiltration_map_eq_atBreak}
\lean{LanglandsFirstMainLemma.norm_unitFiltration_map_inf_eq_criticalNormImage}
\lean{LanglandsFirstMainLemma.norm_unitFiltration_map_sup_eq_belowBreak}
\lean{LanglandsFirstMainLemma.norm_range_inf_unitFiltration_atBreak}
\lean{LanglandsFirstMainLemma.norm_range_sup_unitFiltration_atBreak}
\lean{LanglandsFirstMainLemma.cyclicPrime_differentExponent_eq}
\lean{LanglandsFirstMainLemma.cyclicPrime_trace_lattice_image_eq}
\lean{LanglandsFirstMainLemma.normUnitFiltrationInterval_bijective_belowBreak}
\lean{LanglandsFirstMainLemma.normFieldFiltrationQuotient_bijective_belowBreak}
\lean{LanglandsFirstMainLemma.normUnitFiltrationToBreakSucc}
\lean{LanglandsFirstMainLemma.normUnitFiltrationToBreakSucc_mk}
\lean{LanglandsFirstMainLemma.belowBreakNormTargetImage}
\lean{LanglandsFirstMainLemma.norm_unitFiltration_range_eq_belowBreakNormTargetImage}
\lean{LanglandsFirstMainLemma.normUnitFiltrationToBreakSucc_projection}
\lean{LanglandsFirstMainLemma.belowBreakNormTargetImage_atBreak}
\lean{LanglandsFirstMainLemma.CyclicPrimeNormFiltration}
\lean{LanglandsFirstMainLemma.CyclicPrimeNormFiltration.depth_cases}
\lean{LanglandsFirstMainLemma.CyclicPrimeNormFiltration.herbrand_at_break}
\lean{LanglandsFirstMainLemma.CyclicPrimeNormFiltration.herbrand_first_above}
\lean{LanglandsFirstMainLemma.CyclicPrimeNormFiltration.herbrand_above}
\lean{LanglandsFirstMainLemma.CyclicPrimeNormFiltration.herbrand_denominator_le_next}
\lean{LanglandsFirstMainLemma.CyclicPrimeNormFiltration.herbrand_denominator_lt_next}
\lean{LanglandsFirstMainLemma.CyclicPrimeNormFiltration.herbrand_numerator}
\lean{LanglandsFirstMainLemma.CyclicPrimeNormFiltration.herbrand_denominator}
\lean{LanglandsFirstMainLemma.CyclicPrimeNormFiltration.all_depth_below_bijective}
\lean{LanglandsFirstMainLemma.CyclicPrimeNormFiltration.all_depth_critical_range}
\lean{LanglandsFirstMainLemma.CyclicPrimeNormFiltration.all_depth_critical_kernel_card}
\lean{LanglandsFirstMainLemma.CyclicPrimeNormFiltration.all_depth_critical_cokernel_card}
\lean{LanglandsFirstMainLemma.CyclicPrimeNormFiltration.all_depth_above_bijective}
\lean{LanglandsFirstMainLemma.CyclicPrimeNormFiltration.different}
\lean{LanglandsFirstMainLemma.CyclicPrimeNormFiltration.trace_image}
\lean{LanglandsFirstMainLemma.CyclicPrimeNormFiltration.below_graded_bijective}
\lean{LanglandsFirstMainLemma.CyclicPrimeNormFiltration.below_interval_bijective}
\lean{LanglandsFirstMainLemma.CyclicPrimeNormFiltration.below_field_bijective}
\lean{LanglandsFirstMainLemma.CyclicPrimeNormFiltration.below_full_image}
\lean{LanglandsFirstMainLemma.CyclicPrimeNormFiltration.below_image_inf}
\lean{LanglandsFirstMainLemma.CyclicPrimeNormFiltration.below_image_sup}
\lean{LanglandsFirstMainLemma.CyclicPrimeNormFiltration.critical_kernel}
\lean{LanglandsFirstMainLemma.CyclicPrimeNormFiltration.critical_exact}
\lean{LanglandsFirstMainLemma.CyclicPrimeNormFiltration.critical_kernel_card}
\lean{LanglandsFirstMainLemma.CyclicPrimeNormFiltration.critical_cokernel_card}
\lean{LanglandsFirstMainLemma.CyclicPrimeNormFiltration.critical_full_image}
\lean{LanglandsFirstMainLemma.CyclicPrimeNormFiltration.critical_successor_image}
\lean{LanglandsFirstMainLemma.CyclicPrimeNormFiltration.above_graded_bijective}
\lean{LanglandsFirstMainLemma.CyclicPrimeNormFiltration.above_graded_range}
\lean{LanglandsFirstMainLemma.CyclicPrimeNormFiltration.above_image}
\lean{LanglandsFirstMainLemma.CyclicPrimeNormFiltration.above_successor_image}
\lean{LanglandsFirstMainLemma.CyclicPrimeNormFiltration.field_image_inf}
\lean{LanglandsFirstMainLemma.CyclicPrimeNormFiltration.field_image_sup}
\lean{LanglandsFirstMainLemma.cyclicPrimeNormFiltration}
\leanok
\SourceText{Theorem~\ref{thm:graded-norm} and
Theorem~\ref{thm:norm-filtration}.}
\uses{mod:ramification-different,mod:ramification-traceideals,
mod:ramification-normabovebreak,mod:ramification-normbelowbreak}
\FormalizationContract
The prime-cyclic instance, lower-break witness, residue-degree-one equality,
integral uniformizer, and monogenicity proof remain explicit inputs.
The inverse Herbrand function always converts an upper target depth to a
lower source depth.  The graded denominator is $\Psi(r)+1$, and above the
break it is kept distinct from $\Psi(r+1)$.  Numerator and denominator
containments are separate proofs; all representative formulas are
equalities of quotient classes, and every denominator-deepening square is
explicit.  The all-depth graded map is bridged explicitly to the
below-break, critical, and above-break maps, including the dependent
transport at $\Psi(r)=r$.

The three cases $r<t$, $r=t$, and $t<r$ are exhaustive and disjoint.
Strict below-break splicing excludes the critical graded line, includes the
endpoint $r=t$ only as a trivial interval, and treats
$K^\times/U_K^0$ as the valuation quotient when $t=0$.  Full images below
the break are stated as exact range, intersection, and join equalities, not
as false surjectivity assertions.  The critical cokernel is the target
graded group modulo the actual norm range; lifting from a graded range to
the full critical image uses the separately proved exact
$U_K^{t+1}\to U_F^{t+1}$ endpoint.  Strict above-break statements use the
exact source depths $\Psi(r)$ and $\Psi(r)+1$ and give subgroup equalities,
not inclusions.

The proposition-valued record stores only proved propositions.  Its public
maps, quotient representative formulas, exact-image subgroups, and
projection compatibilities are separate definitions and theorems.  The
assembly uses only modules~\ref{mod:ramification-different},
\ref{mod:ramification-traceideals},
\ref{mod:ramification-normbelowbreak}, and
\ref{mod:ramification-normabovebreak}; it assumes neither
PullbackConductors, NormCharacters, NormRepresentatives, local class field
theory, nor the First Main Lemma.
\end{theorem}
